\section{Motivación}

La visión por computador ha emergido en los últimos años como uno de los campos más relevantes dentro de la inteligencia artificial. Desde el inicio del proyecto queríamos enfrentarnos a un problema en el que esta disciplina tuviera un papel central, pero aplicado a un ámbito que nos resultara cercano y motivador. En nuestro caso, el fútbol era una elección natural, ya que ambos hemos estado vinculados a este deporte y, concretamente, Carlos Mantilla cuenta además con experiencia como entrenador.

Por otro lado, buscábamos desarrollar un proyecto propio que aportara un enfoque diferente respecto a otros trabajos relacionados con el análisis automático de fútbol. La idea inicial era crear una herramienta orientada a dos tipos de usuarios. En primer lugar, personas invidentes o con discapacidad visual que pudieran acudir a un campo de fútbol, grabar el partido con su teléfono móvil y recibir una narración automática de lo que estuviera ocurriendo. En segundo lugar, clubes de fútbol base que retransmiten sus partidos en directo a través de plataformas como YouTube, pero que no disponen de recursos suficientes para contar con comentaristas durante la retransmisión.

Sin embargo, durante el desarrollo del proyecto se identificaron limitaciones importantes asociadas a estos escenarios. La calidad de imagen, el ángulo de grabación, la perspectiva de un teléfono móvil a pie de campo y la cercanía de las cámaras empleadas habitualmente en fútbol base dificultaban de forma significativa tareas como la detección de jugadores, el seguimiento del balón, la estimación de posiciones y la interpretación de acciones. Por este motivo, el enfoque del proyecto evolucionó hacia la narración automática de partidos profesionales, donde las grabaciones broadcast ofrecen mejores condiciones visuales y permiten abordar el problema con mayor fiabilidad.

Esta evolución no elimina la motivación inicial, sino que la reformula como un primer paso hacia un sistema más robusto. Trabajar sobre partidos profesionales permite validar la arquitectura, los modelos y las heurísticas del sistema en un entorno más controlado, dejando como línea futura la adaptación a contextos más complejos y menos estandarizados, como el fútbol base o las grabaciones realizadas desde dispositivos móviles.

Por otro lado, este proyecto supone un reto atractivo e interesante ya que, al ser un problema núnca antes tratado end-to-end, nos resulta emocionante poder abordarlo por primera vez.